\begin{lemma}[Linear identity-holonomy filter]%
    \label{lem:peps_regular_region_holonomy_filter}
    \lean{TNLean.PEPS.exists_regularClosedWalkIdentityFilter,
        TNLean.PEPS.exists_regularClosedWalkGlobalIdentityFilter}
    \leanok
    \uses{thm:peps_regular_region_parent_flatness,
        def:peps_region_parent_interaction}
    Consider a finite graph with regular $G$-injective tensors, a
    connected region $R$, and a closed walk $p$ within $R$. There
    is a physical linear operator $Q$ on $R$ which fixes the
    original regional ground space. On a regional ground vector
    with inserted bonds $u$, it acts as the identity if the
    holonomy around $p$ is the identity, and as zero otherwise.

    Extended by the identity outside $R$, the same operator fixes
    every global vector satisfying the original regional parent
    condition and has the corresponding scalar action on every
    actual inserted closed PEPS. No $G$-isometry is assumed.
    This is a linear operator; Hermiticity and positivity are
    not asserted. It is an auxiliary consequence of the closure
    and accessible-coordinate calculations in
    \cite[Theorem~5.5]{Schuch2010PEPS}.
    These auxiliary regular filters do not establish global closure spanning; see
    \cite{gap:rmp_peps_quantum_double_g_isometry}.
\end{lemma}

\begin{proof}\leanok
    Choose a spanning tree of $R$. The exposed cycle labels of an
    inserted block are simultaneous conjugates of its cycle
    residuals. Evaluation of the closed-walk word commutes with
    simultaneous conjugation. Thus the diagonal detector of
    identity word value acts on each inserted block column by the
    stated scalar. Transport this detector through a local
    $G$-injective inverse and the original physical maps. These
    maps recover every regional ground vector, so the scalar
    action is preserved. An uninserted closed walk has identity
    holonomy. Apply the regional statements to each complementary
    physical slice to obtain the global statements.
\end{proof}
